\section{Stratified Audit Rules and Minimal Dependency Chain}

\subsection{Stratification}

\textbf{Layer 0 (Ontological layer):} A static global state and its algebraic structure are given; no external time or external probability postulates are introduced.\\
\textbf{Layer 1 (Protocol layer):} Scan operators, projection/acceptance windows, and finite-resolution readout kernels are specified; within the protocol, operational time (iteration count) and induced statistics are defined.\\
\textbf{Layer 2 (Identification layer):} Layer 1 structures are mapped to semantic objects such as spacetime/particles/fields; this layer serves only as interpretation and testable interface.

\subsection{Minimal Dependency Chain}

The dependency chain used in this paper is:
\[
\begin{aligned}
&\text{Static lattice ontology}\ \Rightarrow\ \text{Cut-and-project geometry}\ \Rightarrow\ \text{Scan iteration (protocol time)}\\
&\Rightarrow\ \text{Finite-resolution readout}\ \Rightarrow\ \text{Stable sectors and type folding}\ \Rightarrow\ \text{Testable interface propositions}.
\end{aligned}
\]
The key point of this chain is that discreteness and statistics are not externally imposed, but are induced by ``finite-resolution readout.''
